\documentclass[10pt,a4paper]{amsart}
\usepackage[margin=19mm]{geometry}
\usepackage{amsmath,amssymb,mathtools}
\usepackage[scr=rsfs,cal=euler]{mathalfa}
\usepackage{xfrac}
\usepackage[unicode,hidelinks,bookmarks=false]{hyperref}
\newcommand{\ie}{i.e.\ }
\newcommand{\cf}{cf.\ }
\newcommand{\resp}{respectively\ }
\newcommand{\Subsection}{\subsection}
\providecommand{\nobreakdash}{\nobreak\hbox{-}}
\setlength{\emergencystretch}{2em}
\multlinegap0pt
\usepackage{bm}
\usepackage{bbm}
\usepackage{dsfont}
\usepackage{xspace}
\usepackage{cancel}
\usepackage{mathtools}
\usepackage{amsthm}
\makeatletter
\@ifundefined{satz}{\newtheorem{satz}{Satz}}{}
\@ifundefined{corollary}{\newtheorem{corollary}[satz]{Corollary}}{}
\@ifundefined{proposition}{\newtheorem{proposition}[satz]{Proposition}}{}
\makeatother
\title[Numerical Semigroups with $a\_e = 2g+1$]{Numerical Semigroups with $a\_e = 2g+1$}
\author{Michael Hellus, Reinhold H\"ubl, Anton Rechenauer}
\date{Source: arXiv:2604.21948v1 (CC BY 4.0)}
\begin{document}
\maketitle

\noindent\emph{Source and licence.} Numerical Semigroups with $a_e = 2g+1$. Version arXiv:2604.21948v1 (CC BY 4.0), distributed under the Creative Commons Attribution 4.0 International licence, as recorded in the arXiv licence metadata for this exact version and shown on the abstract page https://arxiv.org/abs/2604.21948v1. Full rights notice: licenses/arxiv-2604.21948-CC-BY-4.0.txt.\par\smallskip

\noindent\textit{Editorial scope.} Counted unit: the Proposition whose source label is \texttt{prop:gplusm\_refelcted} in the article (source0.tex lines 316--326, source label \texttt{prop:gplusm\_refelcted}), delivered together with the article's own complete original proof of it (source0.tex lines 328--397, one \texttt{proof} environment of 70 lines). No mathematics is written, repaired or re-derived: statement and proof are the article's own text line for line. Required context retained uncounted: prop:ae\_maximal (lines 203--212); cor:an\_frob (lines 297--301). Every definition, display and label of those lines is retained unchanged. Omitted from this package and not redistributed: 1--202, 213--296, 302--315, 327--327, 398--830. Nothing inside a retained excerpt is omitted or altered: every retained body is delivered exactly as the article's lines are written, so no omission receipt of any kind accompanies this package. Citations: the retained excerpts carry no citation command at all, so no bibliography entry of the article is transcribed and no reference is redistributed. Reference adaptation: none was needed and none was made; every reference inside the delivered unit resolves inside it, so no unresolved reference or citation is produced. Printed slips kept literal: no printed word, symbol, hypothesis or number of the counted statement or of its proof was altered. Layout and class: the article's own class is \texttt{amsart} with packages including fancybox, fancyhdr, epigraph, blindtext, float, amsmath, amssymb, amsthm, eurosym, graphicx, epstopdf, fontenc, inputenc, polynom; the wrapper is the lane's pinned \texttt{amsart} editorial wrapper at 10pt on A4, which restates the theorem-like environments the retained text needs and adds the article's own macro definitions that the retained text needs (none), with extra wrapper packages (bm, bbm, dsfont, xspace, cancel, mathtools). The delivered numbering is the unit's own. Assets: no retained span contains a figure environment, a graphics-inclusion command, a PDF-page inclusion, a table or a TikZ picture; no image file is copied into this package, the package declares no asset dependency and no image file is redistributed with it. Licence: version arXiv:2604.21948v1 of this article is distributed under the Creative Commons Attribution 4.0 International licence, as recorded in the arXiv licence metadata for this exact version and shown on the abstract page https://arxiv.org/abs/2604.21948v1. A full rights notice is delivered at licenses/arxiv-2604.21948-CC-BY-4.0.txt. Characters: every retained excerpt is pure ASCII, so the delivered engine, XeLaTeX with its default Latin Modern text font, reports no missing character.
\begin{proposition}\label{prop:ae_maximal} 
Let $S = \langle a_1, \ldots , a_e \rangle$ be a numerical 
semigroup of genus $g$. Then the following are equivalent: 
\begin{enumerate}
    	\item $a_e = 2g+1$.
    	\item The assignment $\varphi$ given by  $\varphi(n) = 2g+1 - n$ 
	induces a well defined bijection $\varphi: \{1, \ldots, 2g\} \cap S 
	\longrightarrow \{1, \ldots, 2g \} \setminus S$.   
\end{enumerate}
\end{proposition}

\begin{corollary}\label{cor:an_frob}
Let $S = \langle a_1, \ldots, a_e\rangle$ be a numerical semigroup with multiplicity $m$ 
and with $a_e = 2g+1$. Then 
	$$ F(S) = a_e - m = 2g+1 - m $$
\end{corollary}

\begin{proposition}\label{prop:gplusm_refelcted}
Let $S = \langle a_1, \ldots, a_e\rangle$ be a numerical semigroup with multiplicity 
$m=a_1$, genus $g = g(S)$ and with Frobenius element $f = F(S)$. Then the 
following are equivalent: 
\begin{itemize}
\item[i)] $a_e = 2 g+1$. 
\item[ii)] $m+\operatorname{RG}(f, S) = \operatorname{Ap}(S) 
\setminus \{0, f+m\}$. 
\item[iii)] $a_e = f+m$ and $\vert \operatorname{RG}(f, S) \vert = m-2$.
\end{itemize}
\end{proposition}

\begin{proof}
First assume that $a_e = f+m$. Then the following holds 

\begin{enumerate}
\item The map $\iota: \operatorname{RG}(f+m, S) \longrightarrow 
\operatorname{RG}(f, S)$ given by $\iota(L) = L-m$ is well--defined (and injective). 

In fact, if $L \in \operatorname{RG}(f+m, S)$, then $f+m-L \notin S$, hence 
$f+m-L  \leq f$, implying $L \geq m$. As $L$ is a gap, necessarily $L > m$ and 
therefore $L-m \in \{1, \ldots, f-1\}$. Certainly $L-m$ itself is a gap as $L$ is, and so is 
$f-(L-m)=f+m-L$. Hence $L-m \in \operatorname{RG}(f, S)$. 
\item $L \in \operatorname{RG}(f, S) \setminus \operatorname{im}(\iota)$ if and only 
if $L + m \in \operatorname{Ap}(S) \setminus \{0, f+m\}$. 

If $L \in \operatorname{RG}(f, S)$ but $L+m \notin  \operatorname{RG}(f+m, S)$, 
then $L$ is a gap and $f-L = (f+m) - (L+m)$ is a gap of $S$, hence necessarily 
$L+m$ is not a gap of $S$, thus $L + m \in \operatorname{Ap}(S)\setminus \{0, f+m\}$. 

Conversely if $L+m  \in \operatorname{Ap}(S) \setminus \{0, f+m\}$, then 
$L+m \in S$ but $L \notin S$. Furthermore $x=f-L = (f+m) - (L+m) \notin S$, as 
otherwise $f+m = x+(L+m) \in S_+$ contradicting the fact that $a_e = f+m$ is a minimal 
generator. Thus $L \in \operatorname{RG}(f, S) \setminus \operatorname{im}(\iota)$. 
\item $a_e = 2g+1$ if and only if $\operatorname{RG}(f+m, S) = \emptyset$

This is immediate by Proposition~\ref{prop:ae_maximal} as $a_e = f+m$.  
\end{enumerate}
\bigbreak

With this we are able to complete the proof of the proposition: 

Assume i). If $a_e = 2g+1$, then $a_e = f+ m$ by 
Corollary~\ref{cor:an_frob}, hence 
$\operatorname{RG}(f+m, S) = \emptyset$ by (3), and therefore 
	$$ m+\operatorname{RG}(f, S) 
	= \operatorname{Ap}(S) \setminus \{0, f+m\} $$ 
by (2), thus ii) holds.
 
Now assume ii), i.e. assume that $m+\operatorname{RG}(f, S) = 
\operatorname{Ap}(S) \setminus \{0, f+m\}$. 
As $\operatorname{Ap}(S) \setminus \{0, f+m\} $ has $m-2$ 
elements, it remains to show that $a_e = f+m$. 

Assume otherwise. Then $a_e < f+m$ (as any $s > f+m$ is in $S_++S_+$) and 
	$$ f+m = \sum\limits_{i=1}^e n_i \cdot a_i $$
with $n_i \geq 0$ and $\sum\limits_{i=1}^e n_i \geq 2$. As $f \notin S$, 
necessarily $n_1 = 0$, i.e. $ f+m = \sum\limits_{i=2}^e n_i \cdot a_i $, and as 
$a_e < f+m$, 
	$$ a_2, \ldots, a_e \in \operatorname{Ap}(S) \setminus \{0, f+m\} $$ 
Let $i_0 = \operatorname{min}\{i \, \vert \,\, n_i > 0\}$ and write 
	$ f+m = a_{i_0} + s $
with 
	$$ s = \left(n_{i_0}-1\right) \cdot a_{i_0} + \sum\limits_{i=i_0+1}^e n_i \cdot a_i 
	\in S_+ $$
As $a_{i_0} - m \in \operatorname{RG}(f, S)$ 
is a reflected gap of $S$ with respect to $f$,  
	$$ s = f + m - a_{i_0} = f - \left(a_{i_0} - m\right) \quad \notin S $$
a contradiction. Hence $a_e = f+m$ and $\vert \operatorname{RG}(f, S) \vert = m-2$ 
and iii) is true.

If iii) holds, i.e. if $a_e = f+m$ and $\vert \operatorname{RG}(f, S) \vert = m-2$, 
then by (2) we have that 
	$$ m + \operatorname{RG}(f, S) \setminus \operatorname{im}(\iota) 
	= \operatorname{Ap}(S) \setminus \{0, f+m\} $$ 
has $m-2$ elements. Hence $\operatorname{im}(\iota) = \emptyset$, thus 
$\operatorname{RG}(f+m, S) = \emptyset$, and therefore $a_e = 2g+1$ by (3), 
hence we get i). 

This completes the proof of the proposition.  

\end{proof}

\end{document}
